\section{Introduction}
In proof-of-stake Ethereum, validators (proposers) receive on-chain rewards for proposing blocks, but can earn significantly more by exploiting Maximal Extractable Value (MEV) from transaction ordering.  To distribute MEV opportunities, Ethereum adopted Proposer-Builder Separation (PBS) via the MEV-Boost relay system.  In MEV-Boost, validators outsource block construction to off-chain \emph{builders}  through relays who assemble high-value blocks, submit a sealed bid, and then deliver the full block after the proposer commits.  This separation flattens the playing field among validators: proposers no longer need to run sophisticated builder algorithms, and instead accept the highest bid without seeing the transactions.  However, the current PBS ecosystem introduces new trust and centralization issues.

Today, a small number of relays deliver the majority of blocks, creating a central point of failure.  Builders must trust relays not to censor or steal blocks, and proposers must trust that builders will honor bids after the block is committed.  Several recent works and proposals (e.g. EIP-7732) aim to enshrine PBS in-protocol.